\documentclass[10pt,a4paper]{amsart}
\usepackage{amsmath,amssymb,amsthm,graphicx}
% Retained minimal preamble for the Stacks Project chapter more-algebra.tex (More on Algebra).
% The chapter file pinned-original-source.tex carries only \input{preamble} on its first line, so no
% preamble text lives inside the chapter itself.  The Stacks Project uses one single shared preamble for
% every chapter; the snapshot stages it identically as sources/stacks-chow/preamble.tex and
% sources/stacks-etale-cohomology/preamble.tex.  The block below is copied mechanically, line for line,
% from sources/stacks-chow/preamble.tex lines 14-18 (the Xy-pic setup) and 154-217 (the native theorem
% declarations together with the mathematical macros).  Nothing was added, removed or redefined; the
% wrapper class and the amsmath/amssymb/amsthm/graphicx base packages are supplied by the assembler.
% Chapter: more-algebra.tex -- More on Algebra.
% For commutative diagrams we use Xy-pic
\usepackage[all]{xy}

% We use 2cell for 2-commutative diagrams.
\xyoption{2cell}

% Theorem environments.
%
\theoremstyle{plain}
\newtheorem{theorem}[subsection]{Theorem}
\newtheorem{proposition}[subsection]{Proposition}
\newtheorem{lemma}[subsection]{Lemma}

\theoremstyle{definition}
\newtheorem{definition}[subsection]{Definition}
\newtheorem{example}[subsection]{Example}
\newtheorem{exercise}[subsection]{Exercise}
\newtheorem{situation}[subsection]{Situation}

\theoremstyle{remark}
\newtheorem{remark}[subsection]{Remark}
\newtheorem{remarks}[subsection]{Remarks}

\numberwithin{equation}{subsection}

% Macros
%
\def\lim{\mathop{\mathrm{lim}}\nolimits}
\def\colim{\mathop{\mathrm{colim}}\nolimits}
\def\Spec{\mathop{\mathrm{Spec}}}
\def\Hom{\mathop{\mathrm{Hom}}\nolimits}
\def\Ext{\mathop{\mathrm{Ext}}\nolimits}
\def\SheafHom{\mathop{\mathcal{H}\!\mathit{om}}\nolimits}
\def\SheafExt{\mathop{\mathcal{E}\!\mathit{xt}}\nolimits}
\def\Sch{\mathit{Sch}}
\def\Mor{\mathop{\mathrm{Mor}}\nolimits}
\def\Ob{\mathop{\mathrm{Ob}}\nolimits}
\def\Sh{\mathop{\mathit{Sh}}\nolimits}
\def\NL{\mathop{N\!L}\nolimits}
\def\CH{\mathop{\mathrm{CH}}\nolimits}
\def\proetale{{pro\text{-}\acute{e}tale}}
\def\etale{{\acute{e}tale}}
\def\QCoh{\mathit{QCoh}}
\def\Ker{\mathop{\mathrm{Ker}}}
\def\Im{\mathop{\mathrm{Im}}}
\def\Coker{\mathop{\mathrm{Coker}}}
\def\Coim{\mathop{\mathrm{Coim}}}

% Boxtimes
%
\DeclareMathSymbol{\boxtimes}{\mathbin}{AMSa}{"02}

%
% Macros for moduli stacks/spaces
%
\def\QCohstack{\mathcal{QC}\!\mathit{oh}}
\def\Cohstack{\mathcal{C}\!\mathit{oh}}
\def\Spacesstack{\mathcal{S}\!\mathit{paces}}
\def\Quotfunctor{\mathrm{Quot}}
\def\Hilbfunctor{\mathrm{Hilb}}
\def\Curvesstack{\mathcal{C}\!\mathit{urves}}
\def\Polarizedstack{\mathcal{P}\!\mathit{olarized}}
\def\Complexesstack{\mathcal{C}\!\mathit{omplexes}}
% \Pic is the operator that assigns to X its picard group, usage \Pic(X)
% \Picardstack_{X/B} denotes the Picard stack of X over B
% \Picardfunctor_{X/B} denotes the Picard functor of X over B
\def\Pic{\mathop{\mathrm{Pic}}\nolimits}
\def\Picardstack{\mathcal{P}\!\mathit{ic}}
\def\Picardfunctor{\mathrm{Pic}}
\def\Deformationcategory{\mathcal{D}\!\mathit{ef}}

% --- appended 2026-10-12 (Chief): the project own definitions of the reference/slogan/history
% display environments, copied mechanically from the shared Stacks preamble lines 8-12
% (\usepackage{verbatim} plus the three \newenvironment lines, which the chapter preamble keeps
% verbatim from the project preamble). Nothing added, removed or redefined: these environments
% carry no mathematical content and the project itself typesets them as comments. Without these
% three lines a counted statement whose own text uses them fails with Environment slogan undefined.
\usepackage{verbatim}
\newenvironment{reference}{\comment}{\endcomment}
\newenvironment{slogan}{\comment}{\endcomment}
\newenvironment{history}{\comment}{\endcomment}
\usepackage{cleveref}
\crefname{originallemma}{original lemma}{original lemmas}
\Crefname{originallemma}{Original lemma}{Original lemmas}
\makeatletter
\newlabel{more-algebra-lemma-lim-finite-projective-gives-finite-projective}{{0D4B}{source}{}{sourcepointer-1}{}}
\newlabel{more-algebra-lemma-lim-finite-projective-gives-finite-projective@cref}{{[originallemma][0][]0D4B}{[1][0][]source}{}{}{}}
\makeatother
\crefname{originallemma}{original lemma}{original lemmas}
\Crefname{originallemma}{Original lemma}{Original lemmas}
\makeatletter
\newlabel{lemma-limit-henselian}{{0CT7}{source}{}{sourcepointer-2}{}}
\newlabel{lemma-limit-henselian@cref}{{[originallemma][0][]0CT7}{[1][0][]source}{}{}{}}
\makeatother
\crefname{originallemma}{original lemma}{original lemmas}
\Crefname{originallemma}{Original lemma}{Original lemmas}
\makeatletter
\newlabel{lemma-lift-finite-projective-module}{{0D4A}{source}{}{sourcepointer-3}{}}
\newlabel{lemma-lift-finite-projective-module@cref}{{[originallemma][0][]0D4A}{[1][0][]source}{}{}{}}
\makeatother
\crefname{originallemma}{original lemma}{original lemmas}
\Crefname{originallemma}{Original lemma}{Original lemmas}
\makeatletter
\newlabel{lemma-lift-projective}{{0D48}{source}{}{sourcepointer-4}{}}
\newlabel{lemma-lift-projective@cref}{{[originallemma][0][]0D48}{[1][0][]source}{}{}{}}
\makeatother
\crefname{originallemma}{original lemma}{original lemmas}
\Crefname{originallemma}{Original lemma}{Original lemmas}
\makeatletter
\newlabel{lemma-isomorphic-finite-projective-lifts}{{0BC6}{source}{}{sourcepointer-5}{}}
\newlabel{lemma-isomorphic-finite-projective-lifts@cref}{{[originallemma][0][]0BC6}{[1][0][]source}{}{}{}}
\makeatother
\title{The limit of a compatible system of finite flat modules over a surjective system of rings is finite projective}
\author{The Stacks Project Authors}
\date{stacks-project:more-algebra.tex:local-snapshot-20261012; GNU FDL 1.2 or later}
\begin{document}
\maketitle
\setcounter{subsection}{1}
\noindent\textit{Editorial note.} All delivered mathematical text is the author's own, copied from the pinned Stacks Project chapter file more-algebra.tex; no statement or proof step was rewritten, reordered or removed.\par\medskip
\noindent\textit{Editorial note.} The chapter file pinned-original-source.tex carries only \textbackslash{}input\{preamble\} on its first line, so no preamble text lives inside the chapter. The retained file sources/stacks-more-algebra/retained-preamble.tex is a mechanical, line-for-line copy of the project's single shared preamble as staged in the snapshot (Xy-pic setup plus theorem declarations and mathematical macros). Nothing was added, removed or redefined, and the excerpt serves only as preamble, never as a span coordinate.\par\medskip
\noindent\textit{Editorial note.} Both counted spans start at the \textbackslash{}begin line of their environment, so the statement and the author's proof are delivered whole; surrounding prose of the section is deliberately outside the counted spans and is carried as standard prerequisites instead.\par\medskip
\noindent\textit{Editorial note.} The citation fields (source\_id, source\_version, authors, source\_url, license.name, source\_title) are identical to those of every other card this lane has produced from this chapter, so the batch selector treats them as one source.\par\medskip
\noindent\textit{Allowed standard prerequisites (editorial).}\begin{itemize}
\item The inverse limit pair of a system with surjective transition maps and locally nilpotent kernels is henselian. Lemma \textbackslash{}ref\{lemma-limit-henselian\}, cited at proof line 3478. Allows the lifting result for finite projective modules over a henselian pair to be applied.
\item Lifting of finite projective modules over a henselian pair. Lemma \textbackslash{}ref\{lemma-lift-finite-projective-module\}, cited at proof lines 3479-3480. Produces the finite projective module over the limit ring mapping isomorphically to the first stage.
\item Finite flat modules over the stages are finite projective. Lemma \textbackslash{}ref\{lemma-lift-projective\}, cited at proof lines 3493-3494. Makes the lifting result applicable at every stage.
\item Isomorphic finite projective modules lift along the transition maps. Lemma \textbackslash{}ref\{lemma-isomorphic-finite-projective-lifts\}, cited at proof lines 3494-3495. Shows the comparison maps are isomorphisms at every stage.
\end{itemize}\medskip
\ifdefined\hypertarget\hypertarget{sourcepointer-1}{}\fi
\noindent\textit{Source reference.} Original lemma 0D4B: focused-seeds/AGP/sources/stacks-snapshot-20261011/tags.txt, source lines 16446--16446. The lemma's own canonical tag, quoted so that readers can resolve the label inside the delivered text; it points at the counted statement itself and carries no proof content.\par\medskip
\ifdefined\hypertarget\hypertarget{sourcepointer-2}{}\fi
\noindent\textit{Source reference.} Original lemma 0CT7: focused-seeds/AGP/sources/stacks-snapshot-20261011/tags.txt, source lines 16057--16057. Navigation pointer to a statement the counted proof refers to but does not contain; nothing load-bearing is delegated to it.\par\medskip
\ifdefined\hypertarget\hypertarget{sourcepointer-3}{}\fi
\noindent\textit{Source reference.} Original lemma 0D4A: focused-seeds/AGP/sources/stacks-snapshot-20261011/tags.txt, source lines 16445--16445. Navigation pointer to a statement the counted proof refers to but does not contain; nothing load-bearing is delegated to it.\par\medskip
\ifdefined\hypertarget\hypertarget{sourcepointer-4}{}\fi
\noindent\textit{Source reference.} Original lemma 0D48: focused-seeds/AGP/sources/stacks-snapshot-20261011/tags.txt, source lines 16443--16443. Navigation pointer to a statement the counted proof refers to but does not contain; nothing load-bearing is delegated to it.\par\medskip
\ifdefined\hypertarget\hypertarget{sourcepointer-5}{}\fi
\noindent\textit{Source reference.} Original lemma 0BC6: focused-seeds/AGP/sources/stacks-snapshot-20261011/tags.txt, source lines 14269--14269. Navigation pointer to a statement the counted proof refers to but does not contain; nothing load-bearing is delegated to it.\par\medskip
\begin{lemma}
\label{lemma-lim-finite-projective-gives-finite-projective}
Let $A = \lim A_n$ be a limit of an inverse system $(A_n)$ of rings.
Suppose given $A_n$-modules $M_n$ and $A_{n + 1}$-module maps
$M_{n + 1} \to M_n$. Assume
\begin{enumerate}
\item the transition maps $A_{n + 1} \to A_n$ are surjective
with locally nilpotent kernels,
\item $M_1$ is a finite projective $A_1$-module,
\item $M_n$ is a finite flat $A_n$-module, and
\item the maps induce isomorphisms
$M_{n + 1} \otimes_{A_{n + 1}} A_n \to M_n$.
\end{enumerate}
Then $M = \lim M_n$ is a finite projective $A$-module
and $M \otimes_A A_n \to M_n$ is an isomorphism for all $n$.
\end{lemma}

\begin{proof}
By Lemma \ref{lemma-limit-henselian} the pair $(A, \Ker(A \to A_1))$ is
henselian. By Lemma \ref{lemma-lift-finite-projective-module}
we can choose a finite projective $A$-module $P$ and an isomorphism
$P \otimes_A A_1 \to M_1$. Since $P$ is projective, we can successively lift
the $A$-module map $P \to M_1$ to $A$-module maps
$P \to M_2$, $P \to M_3$, and so on. Thus we obtain a map
$$
P \longrightarrow M
$$
Since $P$ is finite projective, we can write $A^{\oplus m} = P \oplus Q$
for some $m \geq 0$ and $A$-module $Q$. Since $A = \lim A_n$ we conclude
that $P = \lim P \otimes_A A_n$. Hence, in order to show that the displayed
$A$-module map is an isomorphism, it suffices to show that the maps
$P \otimes_A A_n \to M_n$ are isomorphisms.
From Lemma \ref{lemma-lift-projective} we see that
$M_n$ is a finite projective module.
By Lemma \ref{lemma-isomorphic-finite-projective-lifts}
the maps $P \otimes_A A_n \to M_n$ are isomorphisms.
\end{proof}

\end{document}
